% Imporditud: tools/import_eksperimendid.py
% Allikas: Eesti füüsikaolümpiaadi eksperimendi ülesannete
%   kogu (Rael Kalda, 2023), ülesanne Ü88, lahendus L88.
\setAuthor{EFO žürii}
\setRound{lõppvoor}
\setYear{2002}
\setNumber{P E2}
\setDifficulty{3}
\setTopic{elektriahelad}
\setVahendid{3 hõõglampi hoidjates, aukudega plastikplaat, 6 neisse aukudesse sobivat vedrukest, 3 metallriba, mida võib kasutada juhtmetena, vooluallikas (3 V) koos ühendusjuhtmetega, multimeeter koos ühendusjuhtmetega}
\setPunktid{12}
\setAllikas{Eksperimendi kogu Ü88, lahendus lk 71}
\setLahendusseis{toores}

\prob{Vooluring hõõglampidega}
Koostage vooluring, milles kolmest ühesugusest hõõglambist kaks on ühendatud omavahel rööbiti ja kolmas lamp selle rööpühendusega jadamisi. Määrake, mitu korda on jadamisi ühendatud lambis eralduv võimsus suurem ühes rööbiti ühendatud lambis eralduvast võimsusest. Kasutatavat multimeetrit võib mõõtepiirkonnal 10 A lugeda ampermeetriks, mis ei mõjuta mõõdetavat voolutugevust (plussklemm kõige ülemine pesa, miinusklemm alumine). Mõõtepiirkondadel 20 V ja 2000 mV võib sedasama mõõteriista lugeda voltmeetriks, mis ei mõjuta mõõdetavat pinget (plussklemm keskmine pesa, miinusklemm alumine). Joonistage töös

kasutatud elektriskeemid.

\hint

\solu
Õigesti visandatud lambid ja ühenduskohad skeemil --- 2 p. Õigesti näidatud voltmeetri ja ampermeetri ühenduskohad skeemil (ampermeeter jadamisi, voltmeeter rööbiti uuritava lambiga) --- 2 p. Ampermeetri ühendamine eraldi mõlemasse rööpharusse pole nõutav. Ülesandes oli ju tehtud eeldus, et lambid on ühesugused, mistõttu kumbagi rööpharu läbib pool sellest voolust, mis kulgeb jadamisi ühendatud lambis. Viimase mõõtmisest ja kahega jagamisest on küllalt. Õigesti määratud pinged --- 2 p. Õigesti määratud voolutugevused --- 2 p. Õigesti arvutatud võimsused (N = IU ) ja võimsuste suhe --- 2 p. Õigesti leitud piirvead --- 2 p.
\probend
